% ------------------------------------------------------------------ 3
\section{Amortized analysis}\label{sec:amortized}

A single \textsc{Extract-Min} can take \(\Theta(n)\) time (after \(n\) inserts the root list has \(n\) trees), and a
single \textsc{Decrease-Key} can cut a whole path. Worst-case bounds per operation are therefore useless for
Dijkstra. What matters is the total cost of a \emph{sequence} of operations, and that is what amortized analysis
bounds. We use the potential method and then restate the result with the accounting and aggregate methods.

\subsection{Cost model and potential function}
Let \(D_0\) be the empty heap and \(D_i\) the heap after the \(i\)-th operation, whose actual cost is \(c_i\). For a
potential \(\Phi\) with \(\Phi(D_0)=0\) and \(\Phi(D_i)\ge0\), the amortized cost is
\[\hat c_i = c_i+\Phi(D_i)-\Phi(D_{i-1}),\qquad\text{so}\qquad
\sum_{i=1}^{N}\hat c_i=\sum_{i=1}^{N}c_i+\Phi(D_N)-\Phi(D_0)\;\ge\;\sum_{i=1}^{N}c_i .\]
The sum of amortized costs is an upper bound on the true total for every sequence.

\begin{definition}[potential]\label{def:phi}
\(\Phi(H)=t(H)+2\,m(H)\), where \(t(H)\) is the number of trees in the root list and \(m(H)\) the number of marked nodes.
\end{definition}

\noindent\(\Phi(D_0)=0\) and \(\Phi\ge0\) always. One unit of cost pays for a constant amount of pointer work. We
charge \textsc{Insert} and \textsc{Union} one unit each; \textsc{Decrease-Key} one unit plus one per cut; and
\textsc{Extract-Min} one unit, plus one per promoted child, plus one per root examined by \textsc{Consolidate}.
Each \textsc{Link} and each iteration of the inner \textbf{while} loop removes one root that was examined, so
there are at most as many of them as examined roots, and the true running time is within a constant factor of this
count. The implementation counts exactly these units (\texttt{FibStats::actual\_cost}).

\subsection{Amortized cost of each operation}

\begin{lemma}[\textsc{Insert}, \textsc{Union}, \textsc{Make-Heap}, \textsc{Minimum}]\label{lem:insert}
\(\hat c_{\textsc{Insert}}=2\), \(\hat c_{\textsc{Union}}=1\), and \(\hat c=1\) for \textsc{Make-Heap} and \textsc{Minimum}.
\end{lemma}
\begin{proof}
\textsc{Insert} adds one unmarked root: \(c=1\), \(\Delta t=1\), \(\Delta m=0\), so \(\hat c=1+1=2\).
\textsc{Union} concatenates two root lists: \(t(H_1\cup H_2)=t(H_1)+t(H_2)\) and the same holds for \(m\), so
\(\Phi\) of the result equals \(\Phi(H_1)+\Phi(H_2)\), \(\Delta\Phi=0\) and \(\hat c=1\).
The other two change nothing.
\end{proof}

\begin{lemma}[\textsc{Extract-Min}]\label{lem:extract}
If the heap has \(n\) nodes and every node has degree at most \(\Dn\), then \(\hat c_{\textsc{Extract-Min}}\le 3\Dn+1\).
\end{lemma}
\begin{proof}
Let \(t=t(H)\) before the call and let \(c'\le\Dn\) be the number of children of the minimum \(z\). Promoting the
children costs \(c'\). After \(z\) is removed the root list has \(t+c'-1\) trees and \textsc{Consolidate} examines each,
so
\[c_i=1+c'+(t+c'-1)=2c'+t .\]
After \textsc{Consolidate} no two roots have the same degree and every degree is in \(\{0,\dots,\Dn\}\), so at most
\(\Dn+1\) roots remain. No node becomes marked (\textsc{Link} and the promotion only clear marks), so
\(\Delta m\le0\) and
\[\hat c_i\le 2c'+t+\big((\Dn+1)-t\big)=2c'+\Dn+1\le3\Dn+1. \qedhere\]
\end{proof}
The term \(t\) cancels: the long root list built up by lazy inserts was paid for in advance by the potential those
inserts left behind.

\begin{lemma}[\textsc{Decrease-Key}]\label{lem:decrease}
\(\hat c_{\textsc{Decrease-Key}}\le5\).
\end{lemma}
\begin{proof}
If no cut happens, \(c_i=1\) and \(\Delta\Phi=0\). Otherwise let \(c\ge1\) be the number of cuts: the cut of \(x\)
and \(c-1\) cascading cuts. Then \(c_i=1+c\). Each cut adds one root, \(\Delta t=c\). Each cascading cut cuts a marked
node and clears its mark, the cut of \(x\) clears its mark if it had one, and at most one node (where the cascade
stops) becomes marked, so \(\Delta m\le 1-(c-1)=2-c\). Hence
\[\hat c_i\le(1+c)+c+2(2-c)=5. \qedhere\]
\end{proof}
This is why marked nodes carry weight two in \(\Phi\): one unit pays for the cut of the node and one for the
new root it creates.

\begin{lemma}[\textsc{Delete}]\(\hat c_{\textsc{Delete}}\le 3\Dn+6\), since it is one \textsc{Decrease-Key} followed by one \textsc{Extract-Min}.\end{lemma}

\subsection{Bounding the maximum degree}
The bound \(\Dn=O(\log n)\) is what makes \cref{lem:extract} logarithmic, and it depends on the marking rule.

\begin{lemma}\label{lem:children}
Let \(x\) be a node and \(y_1,\dots,y_d\) its children in the order in which they were linked to \(x\). Then
\(y_i.\mathit{degree}\ge i-2\) for \(i=1,\dots,d\).
\end{lemma}
\begin{proof}
When \(y_i\) was linked, \(y_1,\dots,y_{i-1}\) were already children of \(x\), so \(x\) had degree at least \(i-1\).
\textsc{Link} joins roots of equal degree, so \(y_i\) also had degree at least \(i-1\). Since then \(y_i\) has lost at
most one child: the first loss marks it, and a second loss would have cut \(y_i\) from \(x\). So its degree is at
least \(i-2\).
\end{proof}

\begin{lemma}\label{lem:size}
A node of degree \(d\) roots a subtree with at least \(F_{d+2}\ge\varphi^{d}\) nodes, where \(F_j\) is the
\(j\)-th Fibonacci number and \(\varphi=(1+\sqrt5)/2\).
\end{lemma}
\begin{proof}
Let \(s_d\) be the smallest possible subtree size of a node of degree \(d\); \(s_0=1\), \(s_1=2\), and \(s_d\) is
non-decreasing in \(d\). By \cref{lem:children}, counting \(x\), its first child and the subtrees of the others,
\(s_d\ge 2+\sum_{i=2}^{d}s_{i-2}\). Induction with \(F_{d+2}=1+\sum_{j=0}^{d}F_j\) gives \(s_d\ge F_{d+2}\), and
\(F_{d+2}\ge\varphi^{d}\) follows from \(\varphi^2=\varphi+1\) by induction.
\end{proof}

\begin{corollary}\label{cor:degree}
Every node of an \(n\)-node Fibonacci heap has degree at most \(\Dn=\lfloor\log_\varphi n\rfloor\le1.4405\log_2 n\).
\end{corollary}
\noindent The array \(A\) in \textsc{Consolidate} is sized by this bound, and the code asserts it on every link.

\subsection{The three methods side by side}
\begin{theorem}\label{thm:sequence}
Starting from an empty heap, any sequence of \(a\) \textsc{Insert}/\textsc{Union} operations, \(b\) \textsc{Extract-Min}
operations, \(c\) \textsc{Decrease-Key} operations and \(d\) \textsc{Delete} operations, in which the heap never holds more
than \(n\) nodes, has total actual cost at most \(2a+b\,(3\Dn+1)+5c+d\,(3\Dn+6)=O(a+c+(b+d)\log n)\).
\end{theorem}
\begin{proof}
Sum the bounds of \cref{lem:insert,lem:extract,lem:decrease} and use \(\sum c_i\le\sum\hat c_i\).
\end{proof}

\paragraph{Accounting method.} The same analysis can be told with credits. Every root holds one credit and every
marked node two. \textsc{Insert} is charged 2: one pays for the work, one is left on the new root. \textsc{Consolidate}
pays for examining and linking a root with that root's credit. \textsc{Decrease-Key} is charged 5: a cascading cut is
paid by the two credits of the marked node it cuts (one for the cut, one placed on the new root), and the charge covers
the first cut, its root credit, and the two credits for the one new mark. The credit balance is exactly
\(\Phi(H)\), so it never goes negative.

\paragraph{Aggregate method for Dijkstra and Prim.} A Dijkstra run on \(n\) vertices and \(m\) edges performs at most
\(n\) \textsc{Insert}, \(n\) \textsc{Extract-Min} and \(m\) \textsc{Decrease-Key} (an edge can lower a label only when its
first endpoint is extracted). By \cref{thm:sequence} the total heap work is at most
\(2n+n(3\Dn+1)+5m=O(m+n\log n)\), the Fredman--Tarjan bound. The heap is emptied at the end, so the next run starts
from \(\Phi=0\) and the bounds add up over the runs of a solver. For Prim on the complete distance network \(D\)
(G1, G3) the counts are \(k\), \(k\) and \(\le k^2\), giving \(O(k^2)\) instead of \(O(k^2\log k)\) with a binary heap.

\subsection{Checking the analysis on real runs}\label{sec:amortized-check}
The proofs above are checked in two ways.

\textbf{Operation by operation.} \texttt{tests/test\_fib\_heap.cpp} runs five random sequences of \(2\times10^5\)
operations (40\% \textsc{Insert}, 30\% \textsc{Decrease-Key}, 25\% \textsc{Extract-Min}, 5\% \textsc{Delete}), compares every result with
\texttt{std::multiset}, verifies the heap invariants every 4096 operations, and asserts after every single operation
that \(\Phi\ge0\) and \(c_i+\Delta\Phi\) does not exceed the bound of its lemma. All \(10^6\) checks pass.
\Cref{fig:amort} plots the running totals: the true cost stays at about \val{amort.test.ratio} of the amortized bound.

\textbf{Inside the solvers.} Each solver reports the heap counters of its run. \Cref{tab:heapcounts} sums them over
all runs of the main benchmark and compares the measured work with the bound of \cref{thm:sequence}.

\begin{figure}[!htbp]\centering
\begin{tikzpicture}
\begin{axis}[width=0.62\linewidth,height=5.2cm,xlabel={operations},ylabel={cumulative cost (units)},
  legend pos=north west,legend style={font=\footnotesize,draw=none},tick label style={font=\footnotesize},
  label style={font=\small},scaled ticks=true,grid=major,grid style={black!8}]
\addplot[cBB,thick,dashed] table[x=op,y=sum_amortized_bound,col sep=comma]{generated/amort_seed1.csv};
\addlegendentry{\(\sum\hat c_i\) (bound)}
\addplot[cG3,thick] table[x=op,y=sum_actual,col sep=comma]{generated/amort_seed1.csv};
\addlegendentry{\(\sum c_i\) (measured)}
\addplot[cG1,thin] table[x=op,y=potential,col sep=comma]{generated/amort_seed1.csv};
\addlegendentry{\(\Phi(D_i)\)}
\end{axis}
\end{tikzpicture}
\caption{Cumulative measured cost against the cumulative amortized bound on one random sequence of
\(2\times10^5\) operations (seed 1). The gap is the slack in the constants of \cref{lem:insert,lem:extract,lem:decrease};
\(\Phi\) stays small because \textsc{Extract-Min} keeps the root list short.}\label{fig:amort}
\end{figure}

\begin{table}[H]\centering\small
\begin{tabular}{@{}lrrrrrrrrr@{}}\toprule
Solver & runs & \textsc{Insert} & \textsc{Extr.-Min} & \textsc{Decr.-Key} & links & cuts & actual & bound & ratio\\\midrule
G1 & 199 & 1.2M & 1.2M & 0.6M & 5.0M & 0.16M & 17.8M & 50.2M & 0.35\\
G2 & 199 & 139.9M & 74.2M & 68.6M & 417.2M & 22.51M & 1380.2M & 3345.5M & 0.41\\
G3 & 199 & 1.3M & 1.3M & 0.9M & 5.3M & 0.26M & 18.9M & 53.9M & 0.35\\
DP & 174 & 1643.5M & 1643.5M & 788.0M & 5957.1M & 165.95M & 21044.4M & 53813.1M & 0.39\\
B\&B & 199 & 1708.3M & 868.2M & 493.8M & 4450.4M & 132.18M & 14744.0M & 60317.1M & 0.24\\
\bottomrule\end{tabular}
\caption{Heap work summed over all runs of the main benchmark (M = million). \emph{actual} is the measured cost in the units of \cref{sec:amortized}; \emph{bound} is \(2\,\#\textsc{Insert}+(3D(N)+1)\,\#\textsc{Extract-Min}+5\,\#\textsc{Decrease-Key}\) from \cref{thm:sequence}, with \(N\) the largest possible heap of the run. For B\&B the counts cover the vertex queue and the open-node queue. On these sparse benchmark graphs there are fewer decrease-keys than extract-mins in every solver; \cref{sec:ablation} shows what this means for running time.}\label{tab:heapcounts}
\end{table}

